\documentclass[10pt,a4paper]{amsart}
\usepackage[margin=19mm]{geometry}
\usepackage{amsmath,amssymb,mathtools}
\usepackage[scr=rsfs,cal=euler]{mathalfa}
\usepackage{xfrac}
\usepackage[unicode,hidelinks,bookmarks=false]{hyperref}
\newcommand{\ie}{i.e.\ }
\newcommand{\cf}{cf.\ }
\newcommand{\resp}{respectively\ }
\newcommand{\Subsection}{\subsection}
\providecommand{\nobreakdash}{\nobreak\hbox{-}}
\setlength{\emergencystretch}{2em}
\multlinegap0pt
\usepackage{amssymb}
\usepackage{mathtools}
\usepackage{bm}
\usepackage{float}
\usepackage{tikz}
\usepackage[shortlabels]{enumitem}
\usepackage{xcolor}
\newtheorem{proposition}{Proposition}
\newcommand{\Ical}{\mathcal{I}}
\newcommand{\Ucal}{\mathcal{U}}
\title{Recursive Filtering and Stochastic Control under Finite Partition-Based Observations}
\author{Saul Díaz-Infante Velasco \and Yofre H. García Gómez \and Jesús Adolfo Minjárez-Sosa}
\date{Source: arXiv:2608.13825v1 (CC BY 4.0)}
\begin{document}
\maketitle

\noindent\emph{Source and licence.} Recursive Filtering and Stochastic Control under Finite Partition-Based Observations. Version arXiv:2608.13825v1 (CC BY 4.0), distributed by arXiv under the Creative Commons Attribution 4.0 International licence, http://creativecommons.org/licenses/by/4.0/, as recorded in the arXiv licence metadata for this exact version and shown on the abstract page https://arxiv.org/abs/2608.13825v1. Full licence notice: \texttt{licenses/arxiv-2608.13825-CC-BY-4.0.txt}.\par\smallskip

\noindent\textit{Editorial scope.} Counted unit: the article's proposition prop:approximate-kernel-consistency (source0.tex lines 2361--2382). The article's complete original proof of it is delivered with it (source0.tex lines 2384--2410, one \texttt{proof} environment of 27 lines). No mathematics is written, repaired or re-derived: the statement and the proof are the article's own lines, in the article's own notation, and no retained line is substituted or re-flowed.  No further source-local context is needed: every cross-reference of every retained line resolves inside the two delivered spans.  Nothing was omitted from any retained span, so the package carries no omission record.  No retained line contains a citation, so the wrapper carries no bibliography and no bibliography entry.  No retained line contains a figure environment, an image-inclusion command or drawing code, so no image is delivered and the package declares no asset dependency.  Editorial formatting only, in addition: an AMS \texttt{amsart} preamble that restates the article's own declarations and macros used by the retained lines, namely the environments \texttt{proposition} on separate counters, together with the standard packages those lines need.  The retained environments print in this document as Proposition 1 in order of appearance, whereas the article prints the same items with the numbering of its own full document.  The complete native source of the article as distributed in the arXiv e-print for this exact version is bundled unchanged as \texttt{sources/207616/source0.tex}.
\begin{proposition}[Kernel consistency on Lipschitz test functions]
	\label{prop:approximate-kernel-consistency}
	For every \(v\in\mathrm{Lip}_b(\Ical)\),
	\begin{equation*}
		\begin{aligned}
			&\sup_{\substack{I\in\mathcal R\\ u\in\Ucal}}
			\Bigg|
			\int_{\mathcal I}
			v(I')\,\mathcal K_N(dI'\mid I,u)
			\\
			&\qquad\qquad
			-
			\int_{\mathcal I}
			v(I')\,\mathcal K(dI'\mid I,u)
			\Bigg|
			\\
			&\qquad\leq
			L_v\varepsilon_N^{\Phi}.
		\end{aligned}
		\label{eq:kernel-consistency-bound}
	\end{equation*}
\end{proposition}

\begin{proof}
	Using the exact branch probabilities in both kernels,
	\[
	\begin{aligned}
		&\left|
		\int v\,d\mathcal K_N(\cdot\mid I,u)
		-
		\int v\,d\mathcal K(\cdot\mid I,u)
		\right|
		\\
		&\quad\leq
		\sum_{k\in\mathsf J}
		\vartheta_k(I,u)
		\left|
		v(k,\Phi_{N,k}(\mu,u))
		-
		v(k,\Phi_k(\mu,u))
		\right|
		\\
		&\quad\leq
		L_v\varepsilon_N^{\Phi}
		\sum_{k\in\mathsf J}\vartheta_k(I,u)
		=
		L_v\varepsilon_N^{\Phi}.
	\end{aligned}
	\]
\end{proof}

\end{document}
